% پاورقی‌های انگلیسی چپ‌به‌راست‌اند، اما شمارۀ آن‌ها با رقم فارسی نمایش داده می‌شود.
\makeatletter
\renewcommand*{\bidi@@LTRfootnotetext@font}{%
    \footnotesize\resetlatinfont
    \def\@makefnmark{\hbox{\@textsuperscript{\xepersian@textdigitfont\@thefnmark}}}%
}
\makeatother

\chapter{مقدمه}

با گسترش اینترنت، کاربران هر روز دیدگاه‌ها، نقدها و تجربه‌های فراوانی دربارۀ محصولات، خدمات، رویدادها و اشخاص منتشر می‌کنند. این متن‌ها منبعی ارزشمند برای شناخت نیازها، ترجیحات و میزان رضایت کاربران به شمار می‌آیند. تحلیل آن‌ها به سازمان‌ها در بهبود محصولات و خدمات، شناسایی نقاط قوت و ضعف و تصمیم‌گیری آگاهانه‌تر کمک می‌کند. بی‌توجهی به این بازخوردها نیز ممکن است بخش مهمی از تجربۀ واقعی کاربران را از فرایند تصمیم‌گیری حذف کند.

حجم زیاد و رشد پیوستۀ متن‌های فضای مجازی بررسی دستی آن‌ها را دشوار و در بسیاری از کاربردها ناممکن کرده است. تحلیل احساسات\LTRfootnote{Sentiment Analysis} یکی از شاخه‌های پردازش زبان طبیعی است که نگرش، ارزیابی یا حالت عاطفی موجود در متن را استخراج می‌کند. این حوزه در تحلیل دیدگاه مشتریان، پایش رضایت از خدمات، خلاصه‌سازی بازخوردها و پشتیبانی از تصمیم‌گیری کاربرد دارد.

در ساده‌ترین صورت، سامانۀ تحلیل احساسات قطبیت متن را در یکی از سه دستۀ مثبت، منفی یا خنثی قرار می‌دهد. بااین‌حال، احساس در همۀ متن‌ها به یک شیوه بیان نمی‌شود. از نظر شیوۀ بیان، نمونه‌های احساسی به دو دستۀ کلی \textbf{احساس صریح} و \textbf{احساس ضمنی} تقسیم می‌شوند.

در احساس صریح، نویسنده نگرش خود را با واژه‌ها یا عبارت‌های ارزیابانۀ آشکار بیان می‌کند. برای مثال، در جملۀ «کیفیت صفحۀ نمایش این لپ‌تاپ عالی است»، واژۀ «عالی» نگرش مثبت نویسنده را نشان می‌دهد. در چنین نمونه‌هایی، یک واژۀ احساسی روشن می‌تواند سرنخ نیرومندی برای تعیین قطبیت باشد.

در احساس ضمنی، نگرش نویسنده بدون یک واژۀ احساسی مستقیم منتقل می‌شود. برای مثال، جملۀ «باتری لپ‌تاپ فقط دو ساعت دوام آورد» هیچ صفتی مانند «ضعیف» یا «نامناسب» ندارد. بااین‌حال، دانش عمومی دربارۀ عمر مطلوب باتری نگرش منفی جمله را آشکار می‌کند. بیان پیامد، مقایسه، توصیه، انتظار برآورده‌نشده یا تجربۀ خاص نیز می‌تواند حامل احساس باشد، بی‌آنکه قطبیت در متن ذکر شود. تشخیص این نوع نگرش موضوع تحلیل احساسات ضمنی\LTRfootnote{Implicit Sentiment Analysis (ISA)} است.

تمایز میان احساس صریح و ضمنی از نظر روش تحلیل اهمیت دارد. روش‌های متکی بر واژگان احساسی در نمونه‌های صریح سرنخ‌های مناسبی می‌یابند، اما در نمونه‌های ضمنی با کمبود شواهد مستقیم روبه‌رو می‌شوند. تحلیل احساسات ضمنی به درک بافت، تشخیص رابطۀ هدف با رویداد و استفاده از دانش عمومی نیاز دارد. سامانه همچنین باید گزارش خنثی را از ارزیابی غیرمستقیم متمایز کند و احساس یک موجودیت را به هدفی دیگر نسبت ندهد.

مدل‌های زبانی بزرگ به دلیل دانش گسترده و توانایی استدلال زبانی، امکان تازه‌ای برای رویارویی با این چالش فراهم کرده‌اند. بااین‌حال، استفاده از این مدل‌ها پرسش‌های مهمی ایجاد می‌کند. آیا استدلال چندمرحله‌ای همواره از پیش‌بینی مستقیم بهتر است؟ اگر مسیر استدلال منحرف شود، چگونه می‌توان خطا را تشخیص داد؟ هنگام اختلاف چند مسیر پیش‌بینی، کدام خروجی باید پاسخ نهایی باشد؟

پژوهش حاضر برای پاسخ به این پرسش‌ها، یک سامانۀ چندمنبعی و مبتنی بر پرامپت\LTRfootnote{Prompt-based} را طراحی و ارزیابی می‌کند. این سامانه پیش‌بینی مستقیم مدل زبانی را در کنار مسیر استدلال چندمرحله‌ای الهام‌گرفته از روش \lr{THOR} قرار می‌دهد. هنگام اختلاف، بازبینی ساختاریافته نوع خطای احتمالی، برچسب پیشنهادی و سطح اطمینان را تولید می‌کند. خروجی کنترل‌گر یک تصمیم میانی است. پیش‌بینی نهایی با انتخاب تنظیم‌شدۀ یکی از منابع مستقیم، استدلالی یا تشخیصی ساخته می‌شود. این فصل مسئله، ضرورت، اهداف، پرسش‌ها، دستاوردها و روش کلی پژوهش را بیان می‌کند.

\section{بیان مسئله}

تمرکز این پژوهش بر تحلیل احساسات ضمنی مبتنی بر هدف است. ورودی مسئله از یک جمله و هدفی مشخص در همان جمله تشکیل می‌شود. هدف می‌تواند محصول، خدمت، ویژگی یا موجودیت مورد ارزیابی باشد. سامانه باید برای هر زوج جمله و هدف، یکی از سه قطبیت مثبت، منفی یا خنثی را پیش‌بینی کند. ازآنجاکه نمونه‌ها ضمنی‌اند، قطبیت از نشانه‌های پراکنده، غیرمستقیم یا وابسته به دانش عمومی استنباط می‌شود.

دشواری نخست تشخیص سرنخ واقعی نظر است. اشاره به یک ویژگی یا رویداد به تنهایی حامل احساس نیست. تفسیر بیش‌ازحد یک گزارۀ خنثی نیز می‌تواند به پیش‌بینی نادرست منجر شود. در مقابل، بعضی پیامدها بدون واژۀ احساسی مطلوب یا نامطلوب‌اند. نادیده‌گرفتن این رابطه باعث ازدست‌رفتن احساس ضمنی می‌شود. دشواری دوم رعایت دامنۀ هدف است. ممکن است جمله چند موجودیت یا چند جنبه داشته باشد. در این حالت، استدلال دربارۀ موجودیتی غیر از هدف تعیین‌شده قطبیت را جابه‌جا می‌کند. ابهام زبانی، نیاز به دانش عرفی و تفاوت دامنه‌ها نیز مسئله را پیچیده‌تر می‌کنند.

یک راه ساده ارائۀ جمله و هدف به مدل زبانی و درخواست مستقیم یکی از سه برچسب است. این راه یک خط پایه\LTRfootnote{Baseline} قوی، کم‌هزینه و روشن فراهم می‌کند. بااین‌حال، مشخص نمی‌کند که مدل بر چه سرنخی تکیه کرده یا در کدام مرحله دچار خطا شده است. راه دیگر تجزیۀ تصمیم به شناسایی جنبه، استخراج نظر ضمنی، استدلال دربارۀ قطبیت و تولید برچسب است. این مسیر ردپای قابل بررسی‌تری ایجاد می‌کند، اما افزایش تعداد گام‌ها دقت را تضمین نمی‌کند. خطای هر گام می‌تواند به گام‌های بعدی منتقل شود و یک استدلال روان نیز به برچسب نادرست برسد.

بررسی خروجی‌های پروژه نشان می‌دهد که هیچ منبعی در همۀ وضعیت‌ها برتری مطلق ندارد. پیش‌بینی مستقیم در بسیاری از نمونه‌ها درست است، درحالی‌که مسیر چندمرحله‌ای در بعضی اختلاف‌ها سرنخ مناسب‌تری فراهم می‌کند. بازبینی تشخیصی نیز گاهی نوع خطا را آشکار می‌کند، اما اعتماد بدون قیدوشرط به آن خطر اصلاح نادرست دارد. ازاین‌رو، مسئلۀ پژوهش فقط تولید یک زنجیرۀ استدلال نیست. سامانه باید چند نشانۀ ناهمگون را کنترل‌شده ترکیب کند و بدون مشاهدۀ برچسب طلایی آزمون، منبع مناسب پاسخ را برگزیند.

برای مطالعۀ این مسئله، داده‌های برچسب‌خوردۀ \lr{SCAPT} مبتنی بر مجموعۀ \lr{SemEval 2014} در دو دامنۀ لپ‌تاپ و رستوران به کار رفته‌اند. پس از استخراج نمونه‌های ضمنی و حذف برچسب‌های نامعتبر، مجموعۀ اصلی شامل ۲۱۸۸ نمونه است. از این تعداد، ۱۷۴۶ نمونه به بخش آموزش و ۴۴۲ نمونه به بخش آزمون تعلق دارند.

ارزیابی اصلی سامانه روی مجموعۀ کامل آزمون با مدل \lr{Qwen3 8B} و از طریق \lr{Ollama} انجام شده است. افزون بر آن، چارچوب آزمایشی با \lr{Gemini 2.5 Flash} نیز روی یک زیرمجموعۀ متوازن ۲۴۰ نمونه‌ای اجرا شده است. این زیرمجموعه ۱۵۰ نمونۀ آموزشی و ۹۰ نمونۀ آزمون دارد. اجرای \lr{Gemini} از طریق یک درگاه سازگار با \lr{OpenAI} انجام شده است. به دلیل تفاوت دامنۀ ارزیابی، نتایج \lr{Gemini} مقایسه‌ای تکمیلی هستند و جایگزین ارزیابی کامل \lr{Qwen} نیستند. وزن هیچ‌یک از مدل‌ها تغییر نکرده است؛ بنابراین، مراحل زبانی صفرنمونه‌ای\LTRfootnote{Zero-shot} و مبتنی بر پرامپت هستند. فقط سیاست انتخاب منبع نهایی با استفاده از بخش آموزش تنظیم می‌شود.

بر این اساس، مسئلۀ مرکزی پایان‌نامه چنین صورت‌بندی می‌شود: چگونه می‌توان پیش‌بینی مستقیم، استدلال چندمرحله‌ای، بازبینی خطا و انتخاب داده‌محور منبع را ادغام کرد؟ این ادغام باید تصمیمی قابل اتکا، قابل تحلیل و بازتولیدپذیر برای تحلیل احساسات ضمنی فراهم کند.

\section{اهمیت و ضرورت پژوهش}

بخش قابل‌توجهی از بازخوردهای واقعی به صورت گزارش تجربه، مقایسه، توصیه، کنایه یا بیان پیامد نوشته می‌شوند. سامانۀ متکی بر واژگان احساسی آشکار ممکن است معنای عملی چنین متن‌هایی را از دست بدهد. این مسئله در تحلیل دیدگاه مشتریان، پایش کیفیت خدمات، خلاصه‌سازی بازخوردها و پشتیبانی از تصمیم‌گیری سازمانی اهمیت دارد. نتیجۀ تحلیل باید نگرش به هدف مشخص را منعکس کند، نه فقط وجود چند واژۀ مثبت یا منفی را.

از منظر علمی، تحلیل احساسات ضمنی چند چالش مهم را دربرمی‌گیرد: استنتاج دانش عمومی، تشخیص هدف، تفکیک توصیف خنثی از ارزیابی و کنترل خطای استدلال. مدل‌های زبانی بزرگ می‌توانند برای این مسائل توضیح و زنجیرۀ استدلال تولید کنند. بااین‌حال، کیفیت ظاهری توضیح صحت پاسخ را تضمین نمی‌کند. ازاین‌رو، پیش‌بینی مستقیم، مسیر استدلال و رفتار انتخاب‌گر نهایی باید هم‌زمان ارزیابی شوند.

ضرورت دیگر پژوهش به فاصلۀ میان یک پرامپت آزمایشی و یک سامانۀ بازتولیدپذیر مربوط است. اجرای قابل اتکا به آماده‌سازی داده، یکنواخت‌سازی خروجی، ادامۀ اجرای قطع‌شده و ذخیرۀ خروجی‌های میانی نیاز دارد. کنترل برچسب‌های نامعتبر، هم‌ترازی ردیف‌ها و تعریف معیارهای ارزیابی نیز ضروری است. معماری پروژه هر مرحله را در خروجی مستقلی ثبت می‌کند تا اثر مؤلفه‌ها و منشأ تصمیم نهایی قابل بررسی باشد.

این پژوهش از نظر روش‌شناختی نیز اهمیت دارد. تنظیم سیاست با دادۀ آموزش و گزارش عملکرد روی همان داده، برداشتی خوش‌بینانه ایجاد می‌کند. ازاین‌رو، معیار اصلی مقایسه، عملکرد روی بخش آزمون است. سیاست انتخاب منبع نیز بدون استفاده از برچسب‌های طلایی آزمون اعمال می‌شود. این تفکیک اعتبار ارزیابی و امکان مقایسۀ منصفانه با خط پایه را حفظ می‌کند.

\section{اهداف پژوهش}

\subsection{هدف اصلی}

هدف اصلی پژوهش طراحی، پیاده‌سازی و ارزیابی یک سامانۀ چندمنبعی برای تحلیل احساسات ضمنی است. این سامانه پیش‌بینی مستقیم را با استدلال چندمرحله‌ای، بازبینی ساختاریافته و انتخاب منبع تنظیم‌شده با دادۀ آموزش ترکیب می‌کند. سامانه بدون ریزتنظیم\LTRfootnote{Fine-tuning} وزن‌های مدل، برای هر زوج جمله و هدف یکی از سه برچسب را تولید می‌کند. خروجی‌های میانی نیز برای تحلیل رفتار و خطا ذخیره می‌شوند.

\subsection{اهداف فرعی}

برای دستیابی به هدف اصلی، اهداف فرعی زیر دنبال می‌شوند:

\begin{itemize}
    \item استخراج و پاک‌سازی نمونه‌های احساس ضمنی از داده‌های \lr{SCAPT} مربوط به دامنه‌های لپ‌تاپ و رستوران و حفظ تقسیم آموزش و آزمون؛
    \item ایجاد یک خط پایۀ پیش‌بینی مستقیم با خروجی سه‌کلاسه و یکنواخت‌شده؛
    \item پیاده‌سازی یک مسیر استدلال چندمرحله‌ای به سبک \lr{THOR} برای تولید جنبه، نظر ضمنی، استدلال قطبیت و برچسب؛
    \item بررسی خودسازگاری\LTRfootnote{Self-consistency} سه‌گانه و رأی اکثریت برای کاهش نوسان مسیر استدلال؛
    \item طراحی بازبینی آگاه از نوع خطا با خروجی ساختاریافته شامل نوع خطا، برچسب پیشنهادی و سطح اطمینان؛
    \item طراحی کنترل‌گری صریح برای مدیریت توافق، اختلاف و اصلاح‌های تشخیصی به عنوان یک تصمیم میانی؛
    \item یادگیری یک سیاست انتخاب منبع از بخش آموزش و اعمال آن بر بخش آزمون، بدون استفاده از برچسب طلایی آزمون؛
    \item مقایسۀ مؤلفه‌ها با معیارهای دقت و امتیاز اف‌یک ماکرو\LTRfootnote{Macro-F۱} و انجام مطالعات حذف مؤلفه و تحلیل خطا؛
    \item اعتبارسنجی هم‌ترازی خروجی‌ها، استخراج نمونه‌های کیفی و فراهم‌سازی اجرای بازتولیدپذیر؛
    \item اجرای چارچوب آزمایشی با \lr{Gemini 2.5 Flash} روی زیرمجموعه‌ای متوازن و مقایسۀ آن با اجرای متناظر \lr{Qwen3 8B}.
\end{itemize}

\section{پرسش‌های پژوهش}

پژوهش حاضر بر پاسخ‌گویی به پرسش‌های زیر متمرکز است:

\begin{enumerate}
    \item پیش‌بینی مستقیم و استدلال چندمرحله‌ای در تحلیل احساسات ضمنی چه تفاوتی دارند؟ آیا استدلال چندمرحله‌ای به تنهایی همواره موجب بهبود می‌شود؟
    \item آیا بازبینی ساختاریافتۀ نوع خطا و یک کنترل‌گر صریح می‌توانند اختلاف میان پیش‌بینی مستقیم و استدلالی را به تصمیمی قابل کنترل‌تر تبدیل کنند؟
    \item آیا انتخاب منبع بر اساس الگوهای آموخته‌شده از دادۀ آموزش می‌تواند در بخش آزمون از اتکای ثابت به یک منبع بهتر عمل کند؟
    \item کدام الگوهای خطا بیشترین نقش را در شکست سامانه دارند؟ آیا ازدست‌رفتن احساس، تفسیر بیش‌ازحد و جابه‌جایی دامنۀ هدف از الگوهای اصلی هستند؟
    \item ظرفیت مدل پایه چه اثری بر مسیر مستقیم، استدلال چندمرحله‌ای و انتخاب منبع دارد؟ مقایسۀ محدود دو مدل چه شواهدی فراهم می‌کند؟
\end{enumerate}

\section{نوآوری‌ها و دستاوردهای پروژه}

نوآوری این پروژه ارائۀ الگوریتمی بنیادین برای مدل‌های زبانی نیست. نوآوری پژوهش در طراحی و ارزیابی ترکیبی منسجم برای تحلیل احساسات ضمنی قرار دارد. دستاوردهای اصلی پروژه عبارت‌اند از:

\begin{itemize}
    \item ایجاد یک خط لولۀ انتهابه‌انتها از تبدیل دادۀ خام تا ارزیابی و تحلیل کیفی، با خروجی‌های میانی قابل ردیابی برای همۀ مراحل اصلی؛
    \item استفاده از \lr{THOR} به عنوان منبع استدلال و نشانۀ مکمل، نه فرض‌کردن آن به عنوان پاسخ نهایی و قطعی سامانه؛
    \item تبدیل بازبینی آزاد به تشخیص ساختاریافتۀ خطا، شامل مجموعه‌ای مشخص از انواع خطا، برچسب اصلاحی و سطح اطمینان؛
    \item تفکیک روشن میان خروجی کنترل‌گر و پیش‌بینی نهایی؛ کنترل‌گر یک تصمیم میانی تولید می‌کند، درحالی‌که پاسخ نهایی از میان منابع مستقیم، \lr{THOR} و تشخیصی انتخاب می‌شود؛
    \item طراحی سیاست انتخاب منبع بر پایۀ الگویی شامل برچسب مستقیم، نوع خطا، اطمینان تشخیصی و دامنه؛ برای هر الگو، بهترین منبع در دادۀ آموزش انتخاب می‌شود و الگوهای دیده‌نشده به منبع مستقیم بازمی‌گردند؛
    \item اعتبارسنجی خودکار زنجیرۀ نهایی از نظر تعداد ردیف‌ها، یکتایی کلیدها، هم‌ترازی خروجی‌ها و اعتبار برچسب‌های نهایی؛
    \item اجرای مقایسۀ تکمیلی \lr{Qwen3 8B} و \lr{Gemini 2.5 Flash} روی یک زیرمجموعۀ متوازن مشترک؛
    \item ارائۀ شواهد عددی و کیفی از اثر ترکیب چندمنبعی: در بخش آزمون ۴۴۲ نمونه‌ای، دقت از \lr{\setpersianfont ۰/۶۷۸۷۳۳} در خط پایۀ مستقیم به \lr{\setpersianfont ۰/۷۲۳۹۸۲} و امتیاز اف‌یک ماکرو از \lr{\setpersianfont ۰/۶۷۴۰۷۵} به \lr{\setpersianfont ۰/۷۱۹۲۰۴} رسیده است. همچنین، سامانۀ نهایی ۲۶ خطای خط پایه را اصلاح کرده و در ۶ مورد، پاسخ درست خط پایه را به پاسخ نادرست تبدیل کرده است.
\end{itemize}

این نتایج به معنای برتری مستقل استدلال چندمرحله‌ای نیستند. در تنظیم آزمایشی حاضر، مسیر \lr{THOR} به تنهایی از خط پایۀ مستقیم ضعیف‌تر بوده است. ارزش این مسیر در فراهم‌کردن استدلال، منبع مکمل و اطلاعات لازم برای بازبینی و انتخاب نهایی آشکار می‌شود. ازاین‌رو، دستاورد محوری پروژه را باید در مدیریت اختلاف و انتخاب آگاهانۀ منبع جست‌وجو کرد.

\section{روش کلی پژوهش}

روش پژوهش از پنج مرحلۀ متوالی و قابل تفکیک تشکیل شده است. در مرحلۀ نخست، فایل‌های برچسب‌خوردۀ \lr{SCAPT} خوانده می‌شوند. برای هر عبارت هدف، جمله، دامنه، تقسیم داده، بازۀ هدف، قطبیت و ضمنی‌بودن استخراج می‌شود. سپس نمونه‌های دارای قطبیت تعارضی حذف و دادۀ سه‌کلاسۀ ضمنی ساخته می‌شود. تقسیم رسمی آموزش و آزمون نیز در تمام مراحل حفظ می‌شود.

در مرحلۀ دوم، دو منبع اصلی پیش‌بینی تولید می‌شوند. در ارزیابی اصلی، هر دو منبع با \lr{Qwen3 8B} و از طریق \lr{Ollama} اجرا می‌شوند. منبع نخست فقط جمله و هدف را دریافت می‌کند و یکی از سه برچسب را برمی‌گرداند. منبع دوم پیاده‌سازی نزدیک به منطق \lr{THOR} است. این منبع تصمیم را به استنتاج جنبه، نظر ضمنی، استدلال قطبیت و تولید برچسب تقسیم می‌کند. مسیر استدلال سه بار با دمای غیرصفر اجرا می‌شود و رأی اکثریت برچسب \lr{THOR} را تعیین می‌کند. هنگام تساوی، نخستین برچسب معتبر میان گزینه‌های مساوی انتخاب می‌شود. یکی از اجراهای هم‌سو با رأی نهایی نیز به عنوان ردپای نماینده نگهداری می‌شود.

در مرحلۀ سوم، فقط هنگامی که پیش‌بینی مستقیم و \lr{THOR} اختلاف دارند، بازبینی آگاه از نوع خطا فعال می‌شود. بازبین با مشاهدۀ جمله، هدف، برچسب مستقیم و ردپای استدلالی، نوع خطا، برچسب پیشنهادی و اطمینان را در قالبی ساختاریافته تولید می‌کند. انواع خطا شامل احساس ازدست‌رفته، تفسیر بیش‌ازحد خنثی، جابه‌جایی دامنۀ هدف و ناهماهنگی جنبه و نظر هستند. ناسازگاری استدلال و برچسب و نبود شواهد کافی نیز در این مجموعه قرار دارند. سپس یک کنترل‌گر قاعده‌محور توافق، قابلیت اصلاح و سطح اطمینان را بررسی می‌کند و تصمیمی میانی می‌سازد.

در مرحلۀ چهارم، انتخاب منبع نهایی انجام می‌شود. نمونه‌های بخش آموزش بر اساس چهار ویژگی برچسب مستقیم، نوع خطا، اطمینان تشخیصی و دامنه گروه‌بندی می‌شوند. در هر گروه، تعداد پیش‌بینی‌های درست سه منبع مستقیم، \lr{THOR} و تشخیصی محاسبه می‌شود و منبع دارای امتیاز بیشتر به آن الگو اختصاص می‌یابد. این جدول تصمیم روی بخش آزمون اعمال می‌شود. اگر الگویی در آموزش دیده نشده باشد، منبع مستقیم به عنوان گزینۀ پیش‌فرض انتخاب می‌شود. در این مرحله هیچ برچسب طلایی از بخش آزمون در دسترس انتخاب‌گر نیست.

در مرحلۀ پنجم، همۀ روش‌ها با دقت و امتیاز اف‌یک ماکرو ارزیابی می‌شوند. امتیاز اف‌یک ماکرو در کنار دقت گزارش می‌شود تا عملکرد هر سه کلاس در نتیجه نقش یکسان‌تری داشته باشد. نتیجۀ اصلی از اجرای \lr{Qwen3 8B} روی مجموعۀ کامل آزمون به دست می‌آید. برای مقایسۀ ظرفیت مدل، همان چارچوب با \lr{Gemini 2.5 Flash} روی زیرمجموعۀ متوازن ۲۴۰ نمونه‌ای اجرا می‌شود. اجرای متناظر \lr{Qwen3 8B} روی همین زیرمجموعه مبنای مقایسۀ منصفانه با \lr{Gemini} است. نتایج این مقایسه به دلیل محدودۀ متفاوت، جایگزین نتیجۀ اصلی نیستند.

در تمام اجراهای گزارش‌شده، وزن مدل‌های زبانی ثابت می‌ماند و هیچ ریزتنظیمی انجام نمی‌شود. افزون بر معیارهای عددی، تحلیل خطا، نمونه‌های کیفی، مطالعات سیاست انتخاب و هم‌ترازی فایل‌ها نیز بررسی می‌شوند.

\section{ساختار پایان‌نامه}

این پایان‌نامه در شش فصل تنظیم شده است:

\begin{enumerate}
    \item فصل نخست با عنوان «مقدمه» مسئله، ضرورت، اهداف، پرسش‌های پژوهش، دستاوردها و روش کلی را معرفی می‌کند.
    \item فصل دوم با عنوان «تعاریف اولیه» مفاهیم پایۀ تحلیل احساسات، مدل‌های زبانی، استدلال زنجیره‌ای و معیارهای ارزیابی را تعریف می‌کند.
    \item فصل سوم با عنوان «ادبیات تحقیق» پژوهش‌های مرتبط با تحلیل احساسات ضمنی، روش \lr{THOR}، خوداصلاحی و انتخاب منبع را مرور می‌کند.
    \item فصل چهارم با عنوان «کار پیشنهادی» معماری سامانه، آماده‌سازی داده، مسیرهای پیش‌بینی، بازبینی و سیاست انتخاب منبع را تشریح می‌کند.
    \item فصل پنجم با عنوان «نتایج آزمایش‌ها» تنظیمات آزمایش، نتایج کمی، مطالعات حذف مؤلفه، تحلیل کیفی و مقایسۀ مدل‌ها را ارائه می‌دهد.
    \item فصل ششم با عنوان «نتیجه‌گیری و کارهای آتی» یافته‌ها را جمع‌بندی می‌کند، به پرسش‌های پژوهش پاسخ می‌دهد و کارهای آینده را شرح می‌دهد.
\end{enumerate}
